% Source: https://artofproblemsolving.com/community/c3h2552648_math_puzzle
% Stable identity: aops-topic-2552648-q01
The argument treats $J$ as a real number that can be subtracted from both sides. But an infinite sum only has a value if its partial sums approach a limit, and here the partial sums are
\[1+2+\cdots+2^{n-1}=2^n-1,\]
which grow without bound. So the series diverges and $J$ is not a number; the step "subtract $2J$" is meaningless. (Positive terms could never add up to a negative number anyway: every partial sum is at least $1$.)

The same manipulation is valid for $1+r+r^2+\cdots$ with $|r|<1$, because there the partial sums $\dfrac{1-r^n}{1-r}$ do converge, to $\dfrac1{1-r}$. Formally substituting $r=2$ into $\dfrac{1}{1-r}$ is exactly what produces the false "$-1$".
